% A-only report fragment. Include where \Owner and \MemberFour are defined.
\section{Basic Analysis}
\label{sec:basic}
\Owner{\MemberFour}{0.75}

\subsection{Idea}
A is a bounded concrete testing analysis for line reachability. It searches
for an input that executes the target instruction, first using common small
and boundary values, then fresh random inputs. Unlike S, A evaluates branch
conditions during execution, but does not use them to guide input generation.
A target hit provides a concrete witness of reachability. A therefore returns
only \texttt{reachable} or \texttt{unknown}, never \texttt{unreachable}.
Although grouped with the simplest analysis rungs, A is dynamic rather than
a purely syntactic static analysis.

\subsection{How it works}
\paragraph{A: Basic}
A first tries the fixed seeds $0,1,-1,2,-2,m,m+1,M-1,M$, where
$m=-2^{31}$ and $M=2^{31}-1$ are the Java \texttt{int} boundaries.
For one argument this gives nine inputs; for two arguments it gives their
Cartesian product of 81 pairs. These combinations include equal arguments,
different argument orders and boundary values.

If no fixed input reaches the target, A tries up to 1,000 additional random
inputs per query. Each argument is sampled independently and uniformly from
the full signed 32-bit integer domain, with replacement. Each execution starts
with fresh locals, operand stack and program counter. The first target hit
stops the search and returns \texttt{reachable}, together with the input as
a witness. If the budget is exhausted, A returns \texttt{unknown}.
A fresh random-generator seed is chosen for each run and recorded for replay
with the same artifacts, query order, budget and Python environment.

For the Group 15 dataset, we reuse the existing \texttt{javap} loader and
concrete bytecode interpreter. The runner consumes the validator's answer-free
query manifest and checks source/class hashes, method descriptors and target
mappings. It observes entry into the designated instruction, not merely the
method's return value. Expected labels, supplied witnesses, validator
representatives and program constants are not used to select inputs.
Invalid mappings and unsupported bytecode cause errors rather than verdicts.
This is a standalone integration, not a verified integration with the team's
JPAMB interpreter.

\subsection{Guarantees}
If A returns \texttt{reachable}, it has observed a concrete input whose
execution enters the target instruction. With correct interpreter semantics
and target mapping, this input is an existential witness:
\[
    A(P,t)=\texttt{reachable}
    \;\Longrightarrow\;
    \exists x \in D:\ P(x)\mbox{ executes }t,
\]
where $D$ is the permitted input domain, including argument pairs for
two-argument methods. A is therefore sound for its reachability claims,
but incomplete: a reachable target may remain undiscovered.
Exhausting a finite testing budget does not exclude other successful inputs,
so A never returns \texttt{unreachable}. These guarantees concern the supported
static, loop-free integer methods without calls or exception handlers.

\subsection{What it cannot do}
A does not solve path conditions or adapt its input choices to guards.
For example, the equality $x=100003$ is not covered by the fixed seeds.
A single uniform random draw hits that value with probability $1/2^{32}$,
so a small random budget can easily miss it. Similarly, failed executions
cannot establish that contradictory conditions make a target unreachable.

On all 30 queries, using JDK 17.0.21, Python 3.12.10, generator seed 0 and
the default budget, A returned nine \texttt{reachable} and 21 \texttt{unknown}.
All nine witnesses came from fixed seeds and were replayed successfully on
the JVM. The unknown results include all 15 unreachable targets and six
reachable targets requiring exact values absent from the fixed seeds.
Random runs may differ. These observations concern this constructed dataset,
not arbitrary Java programs, and motivate later rungs that exploit program
information instead of choosing inputs blindly.
